\chapter[title={Stochastik},reference={Stochastik}]
\section[title={2024-08-28 - Einleitung},reference={2024-08-28-Einleitung}]

\startframedexample
{\bf Statistik vs Stochastik}

Stochastik ist die Vorhersage

Statistik ist die Auswertung der Vargangenheit

\stopframedexample

\startframedtask
{\bf Satz:} Die Wahrscheinlichkeiten der Egebnisse eines
Zufallsexperiments sind Zahlen im intervall $\[0;1\]$ mit Summe $1$. Sie
bilden eine {\em Wahrscheinlichkeitsverteilung}. Sie sind die Prognosen
für die relativen Häufigkeiten bei vielen Versuchswiederholungen.

\stopframedtask

\startframedtask
{\bf Definition:} Wenn jedem Ergebnis eines Zufallsexperiments ein
Zahlenwert zugeordnet wird, spricht man von einer {\bf Zufallsgröße}.
Die {\bf Wahrscheinlichkeitsverteilung} ener Zufallsrgöße $X$ ist eine
Tabelle, bei der jedem Wert $k$ von $X$ die Wahrscheinlichkeit $P(X=k)$
zugeordnet ist. Für eine Zufallsgröße $X$ mit den Werten
$x_1,x_2, ..., x_n$ heißt $\mu = x_1 \cdot
P(X=x_1) + x_2 \cdot P(X = x_2)... + x_n \cdot P(X = x_n)$ {\bf
Erwartungswert} von $X$. Er gibt an, welchen Mittelwert man bei
ausreichend großer Versuchsanzahl auf lange Sicht erwartet.

\stopframedtask

\section[title={Beipsiel: Faires
Spiel},reference={beipsiel-faires-spiel}]

\margintext{Als {\bf fair} bezeichnet man ein Spiel, bei dem der Erwartungswert für
den Gewinn null ist.

Gewinn = Auszahlung - Einsatz
}

Beim Glücksspeil mit einem Würfel soll das Doppelte der Augenzahl (in
Euro) ausgezahlt werden.

\startitemize[a][stopper=)]
\item
  Bestimmen Sie die Auszahlung, die der Spieler im Mittel erwarten kann.
\item
  Geben Sie an, wie hoch der Einsatz sein muss, damit das Glücksspiel
  fair ist.
\stopitemize

\textpartdivider

\startitemize[a][stopper=)]
\item
  Wegen $\mu = \frac{1}{6}(2+4+6+8+10+12)=\frac{42}{6}=7$
\item
  Dei Einsatz sollte dem Erwartungswert entsprechen. So hat der anbieter
  des Glückspieles zwar keinen gewinn, aber auf lange sicht auch keinen
  direkten Verlust und die Teilnehmenden habe eine faiere Chanche.
\stopitemize



\section[title={Aufgaben:},reference={aufgaben}]

\section[title={S. 238 Aufgabe
1},reference={s.-238-aufgabe-1}]

\startplacetable[location=none]
\startxtable
\startxtablehead[head]
\startxrowgroup[lastrow]
\startxrow
\startxcell Gewinn (Chips) \stopxcell
\startxcell Wahrscheinlichkeit \stopxcell
\stopxrow
\stopxrowgroup
\stopxtablehead
\startxtablebody[body]
\startxrow
\startxcell -1 \stopxcell
\startxcell $\frac{1}{4}$ \stopxcell
\stopxrow
\startxrow
\startxcell 0 \stopxcell
\startxcell $\frac{1}{2}$ \stopxcell
\stopxrow
\startxrowgroup[lastrow]
\startxrow
\startxcell 1 \stopxcell
\startxcell $\frac{1}{4}$ \stopxcell
\stopxrow
\stopxrowgroup
\stopxtablebody
\startxtablefoot[foot]
\stopxtablefoot
\stopxtable
\stopplacetable

\startitemize[a,packed][stopper=)]
\item
  Warum ist die Tabelle korrekt?
\stopitemize

Es gibt vier verschiedene Möglichkeiten für die Münzen zu fallen wenn
$0$ gleich Zahl und $1$ gleich Kopf, sind das folgende Möglichkeiten:
$00,01,10,11$. Jede dieser Möglichkeiten tritt mit der selben
Wahrscheinlichkeit auf ($\frac{1}{4}$). In zwei der Fällen ($10,01$)
bekommt man einen Chip zurück, ist also selber auf $0$. Wenn der Fall
$00$ auftritt, verliert man den gesetzten Chip und bei $11$ gewinnt man
einen

\textpartdivider

Roulett

Erwartungswert bei 1Euro Einsatz:

\startformula 
\mu = \frac{1}{37}(36) = \frac{36}{37} \approx 0.97
 \stopformula
